% Ad Familiares 10.31 (ad-familiares-10-31) --- English translation
% Translated by Claude (Anthropic) from the Latin (Perseus).
% Style guide: STYLE.md.

\ciceroLetterOpener{C. ASINIVS POLLIO CICERONI S. D.}

\ciceroSection{1} It should seem to you not in the least surprising that I have written nothing about the state of the commonwealth since the resort to arms. For the pass of Castulo, which has always held up our couriers --- though now it is rendered the more dangerous by the increase of banditry --- still is nothing like the hazard it once was, compared with men stationed at every point on both sides to search out couriers and detain them. And so unless a letter came by ship, I should not know at all what was happening over there. As it is, now that the seas are open for sailing, I shall, with my opportunity, write to you as eagerly and as often as I can.

\ciceroSection{2} That I should be swayed by the talk of the man whom, though no one wants to look at him, men nonetheless do not hate as he deserves --- there is no danger of that; for he is so hateful to me that I think nothing not bitter that I share with him. My nature and my studies in any case draw me toward a passion for peace and liberty. Hence I often lamented the very beginning of the civil war; and when it was no longer permitted me to belong to neither side, since on both I had powerful enemies, I shunned the camp in which I plainly knew I should not be safe from the plots of my enemy; driven where I least wished to go, lest I be left in the lurch, I went into danger plainly and without hesitation.

\ciceroSection{3} Caesar, however, for treating me --- newly known to him, in the midst of such fortune --- as one of his oldest intimates, I loved with the deepest devotion and loyalty. What I was free to do by my own judgement, I did in such a way that the best men gave it their fullest approval; what I was ordered to do, I did at such a time and in such a manner that it was plain it had been imposed on an unwilling man. The very unjust odium attaching to that conduct was capable of teaching me how pleasant liberty is and how miserable life is under a domination. So if what is now afoot is that all things should be again in the power of one man --- whoever he is --- I declare myself his enemy, and there is no danger whatever that I shall shirk or refuse on liberty's behalf.

\ciceroSection{4} But the consuls had given me no instruction --- neither by a senatorial decree nor by a letter of their own --- as to what I was to do; for I received only one letter, after the Ides of March at last, from Pansa, in which he urges me to write to the Senate that I and my army shall be in its power. Which, since Lepidus was haranguing in public and writing to everyone that he was in agreement with Antony, was utterly out of the question; for with what supplies, against his will, was I to lead my legions through his province? Or if I had got past everything else, could I really have flown across the Alps, which are held by his garrison? Add to this that letters could not be carried through on any terms; for they are searched in six hundred places, and then also the couriers are detained by Lepidus.

\ciceroSection{5} That at Corduba I said before the assembly that I would hand over the province to no one except a man sent by the Senate --- no one will call into doubt. As to what struggles I had about handing over the Thirtieth Legion, what need have I to write? Once it had been handed over, who can fail to know how much weaker I should have been for the commonwealth's purposes? For believe nothing fiercer or more pugnacious than this legion. So take me for a man, in the first place, set most eagerly on peace (for I am plainly intent on having every citizen safe), and, in the second, ready to vindicate myself and the commonwealth into liberty.

\ciceroSection{6} That you count an intimate of mine in the number of your own people is more agreeable to me than you suppose; I envy him, even so, for walking and joking with you. You will ask how dearly I prize it. If ever it is permitted me to live in leisure, you shall find out; for I shall not depart a footstep's breadth from your side. The one thing I greatly wonder at is that you have not written to me whether I can do more for the commonwealth by staying in my province or by leading the army into Italy. As for me, though it is safer and less laborious to stay, still, because I see that at such a time legions are wanted far more than provinces --- particularly ones that can be recovered with no trouble --- I have determined, as matters stand, to set out with the army. You will learn everything from the letter I have sent to Pansa, of which I have sent you a copy. 16 March, Corduba.
